\documentclass[11pt]{article}
\usepackage[T1]{fontenc}
\usepackage{lmodern}
\usepackage[a4paper,margin=28mm]{geometry}
\usepackage{amsmath,amssymb,amsthm}
\usepackage{microtype}
\usepackage[colorlinks=true,linkcolor=blue,citecolor=blue,urlcolor=blue]{hyperref}
\hypersetup{pdftitle={A counterexample to the facet bound for reflexive polytopes},pdfauthor={Kaiwen Jin}}
\newtheorem{theorem}{Theorem}
\newtheorem{lemma}[theorem]{Lemma}
\newtheorem{proposition}[theorem]{Proposition}
\newcommand{\conv}{\operatorname{conv}}
\newcommand{\inte}{\operatorname{int}}
\newcommand{\vertices}{\operatorname{Vert}}
\newcommand{\supp}{\operatorname{supp}}
\title{A counterexample to the facet bound\newline for reflexive polytopes}
\author{Kaiwen Jin\\[2pt]\small\href{mailto:yui.ui.siki@gmail.com}{yui.ui.siki@gmail.com}}
\date{}
\begin{document}
\maketitle

\begin{abstract}
A lattice polytope is reflexive if it contains the origin in its interior
and its polar is also a lattice polytope.
We construct a $45$-dimensional reflexive lattice polytope with
$2^{32}\binom{40}{8}+5$ facets, exceeding $6^{45/2}$.
This disproves the unrestricted facet bound conjectured by Nill.
The construction starts with a lattice polytope having many vertices and
a polar whose denominators divide $32$. Five successive pyramids multiply
the original polar coordinates by $32$ while adding only five vertices.
Taking the polar gives the required counterexample.
\end{abstract}

\section{The facet bound}

A lattice polytope in $\mathbb R^d$ is the convex hull of finitely many
points of $\mathbb Z^d$. A full-dimensional lattice polytope $P$ is
\emph{reflexive} if $0\in\inte P$ and its polar
\[
 P^*=\{y\in\mathbb R^d:\langle x,y\rangle\leq1
                         \text{ for every }x\in P\}
\]
is also a lattice polytope. Here interior always means ordinary Euclidean
interior. Reflexive polytopes occur in the study of Gorenstein toric Fano
varieties; their face numbers describe part of the associated combinatorial
complexity \cite{Nill}.

Nill conjectured that a $d$-dimensional reflexive polytope has at most
$6^{d/2}$ vertices \cite[Conjecture~5.2]{Nill}. Polarity interchanges vertices
and facets and preserves reflexivity, so this is equivalent to
\begin{equation}\label{eq:bound}
 f_{d-1}(P)\leq6^{d/2},
\end{equation}
where $f_{d-1}(P)$ is the number of facets of $P$.
Mori, Mori and Ohsugi restate this facet conjecture and prove the bound for
twinned chain polytopes \cite[Conjecture~1.1 and Theorem~1.2]{MMO}.
The unrestricted statement is false.

\begin{theorem}\label{thm:main}
There exists a $45$-dimensional reflexive lattice polytope $P$ such that
\[
 f_{44}(P)=330303106984181765>6^{45/2}.
\]
\end{theorem}

The construction uses a cube cut by a bound on the sum of the absolute
coordinates. It has many integral vertices, but its polar is initially
rational. A pyramid with base at height $1$ and apex at height $-1$ doubles
the old polar coordinates. Repeating this operation clears the denominators
at a cost of only one dimension and one vertex per step.
The general idea of realizing lattice polytopes as faces of reflexive
polytopes predates this construction: Haase and Melnikov prove that every
lattice polytope has such an embedding \cite[Proposition~2.2]{HM}.
Here the explicit dimension and vertex count give the strict violation
of \eqref{eq:bound}. We do not determine the smallest dimension in which
that bound fails.

\section{A clipped cube and its pyramids}

For a polytope $K$, write $\vertices(K)$ for its vertex set and
$h_K(y)=\max_{x\in K}\langle x,y\rangle$ for its support function.
We use the elementary identity $K^{**}=K$ for a full-dimensional
polytope containing $0$ in its interior. One direct justification is as
follows. If $z\notin K$, choose a nearest point $q\in K$ and set $a=z-q$.
Minimality along every segment from $q$ to $x\in K$ gives
$\langle a,x-q\rangle\leq0$. Thus
$h_K(a)=\langle a,q\rangle<\langle a,z\rangle$.
An interior ball implies $h_K(a)>0$, so $a/h_K(a)\in K^*$ separates $z$
from $K^{**}$. The other inclusion follows from the definition.
An interior ball in $K$ bounds $K^*$, and boundedness of $K$ supplies an
interior ball in $K^*$; in particular both polars have full dimension.

For integers $m\geq2$ and $1\leq r<m$, define
\begin{equation}\label{eq:base}
 R_{m,r}=\left\{x\in\mathbb R^m:
           |x_i|\leq1\ (1\leq i\leq m),\quad
           \sum_{i=1}^m|x_i|\leq r\right\}.
\end{equation}
Let $e_i$ denote the $i$th coordinate unit vector, and write
$\supp v=\{i:v_i\neq0\}$ for $v\in\mathbb R^m$.

\begin{lemma}\label{lem:base}
The polytope $R_{m,r}$ has $0$ in its interior and
\begin{align}
 \vertices(R_{m,r})&=\{v\in\{-1,0,1\}^m:|\supp v|=r\},
                                                    \label{eq:basevertices}\\
 |\vertices(R_{m,r})|&=2^r\binom mr,                     \label{eq:basecount}\\
 R_{m,r}^*&=\conv\left(\{\pm e_i:1\leq i\leq m\}
                    \cup\{s/r:s\in\{-1,1\}^m\}\right).
                                                    \label{eq:basepolar}
\end{align}
\end{lemma}

\begin{proof}
The last constraint in \eqref{eq:base} is equivalent to
$\langle s,x\rangle\leq r$ for every $s\in\{-1,1\}^m$.
Hence $R_{m,r}$ is a bounded intersection of finitely many halfspaces.
All inequalities are strict in a neighborhood of $0$.

In any fixed orthant, set $z_i=|x_i|$. The constraints become
$0\leq z_i\leq1$ and $\sum_i z_i\leq r$.
If the sum constraint is inactive, a fractional coordinate can be perturbed
in both directions. If it is active and two coordinates are fractional,
they can be perturbed in opposite directions. Exactly one fractional
coordinate is impossible when the sum is the integer $r$ and the other
coordinates are $0$ or $1$. Thus every vertex of the orthant section has
only $0$ and $1$ coordinates. A vertex of the whole polytope must be a
vertex of each orthant section containing it, so it is ternary with at most
$r$ nonzero coordinates.

If a ternary vector has fewer than $r$ nonzero coordinates, choose a zero
coordinate $j$. The vector is the midpoint of $x+e_j$ and $x-e_j$, both
in $R_{m,r}$, and therefore is not a vertex. Conversely, a ternary vector
$v$ with exactly $r$ nonzero coordinates uniquely maximizes
$\langle v,x\rangle$: equality with the maximum $r$ fixes all supported
coordinates, and the sum bound forces the remaining coordinates to vanish.
This proves \eqref{eq:basevertices}; choosing the support and its signs
gives \eqref{eq:basecount}.

Let $K$ be the convex hull on the right of \eqref{eq:basepolar}.
It contains $\conv\{\pm e_i\}$ and thus has $0$ in its interior.
Its polar inequalities are exactly \eqref{eq:base}, so $K^*=R_{m,r}$.
Bipolarity gives \eqref{eq:basepolar}. The displayed list is a generating
set; its points need not all be vertices.
\end{proof}

For any full-dimensional polytope $R\subset\mathbb R^a$ containing $0$ in
its interior, form the pyramid
\begin{equation}\label{eq:pyramid}
 T(R)=\conv\big((R\times\{1\})\cup\{(0,-1)\}\big).
\end{equation}

\begin{lemma}\label{lem:pyramid}
The polytope $T(R)$ has dimension $a+1$, contains $0$ in its interior,
and has the vertices $(v,1)$ for $v\in\vertices(R)$ together with $(0,-1)$.
Moreover,
\begin{equation}\label{eq:pyramidpolar}
 T(R)^*=\conv\big((2R^*\times\{-1\})\cup\{(0,1)\}\big).
\end{equation}
\end{lemma}

\begin{proof}
The base is an $a$-dimensional exposed face, and the apex lies outside its
affine hull. At height $t\in(-1,1)$ the section is $\frac{1+t}{2}R$.
An interior ball in $R$ therefore supplies a neighborhood of $(0,0)$
in $T(R)$. The base vertices remain vertices, since any decomposition at
height $1$ stays in the base. The apex uniquely minimizes the last
coordinate, and a finite convex hull has no vertices outside its generators.

Testing the base and apex in the polar definition gives
\[
 (y,t)\in T(R)^*\quad\Longleftrightarrow\quad
 t\geq-1,\qquad h_R(y)+t\leq1.
\]
Since $h_R(y)\geq0$, one also has $t\leq1$. For $t<1$ set
$\lambda=(1-t)/2$ and $u=y/(1-t)$. Then $u\in R^*$ and
\[
 (y,t)=\lambda(2u,-1)+(1-\lambda)(0,1),\qquad 0<\lambda\leq1.
\]
At $t=1$, the inequality $h_R(y)\leq0$ and an interior ball force $y=0$.
Conversely, every $(2u,-1)$ with $u\in R^*$ and the point $(0,1)$ satisfy
the polar inequalities. This proves \eqref{eq:pyramidpolar}.
\end{proof}

\section{An integral polar and the counterexample}

Fix integers $k\geq1$ and $m>2^k$, put $r=2^k$, and define
$Q_0=R_{m,r}$ and $Q_j=T(Q_{j-1})$ for $1\leq j\leq k$.
The base polar in \eqref{eq:basepolar} has denominators dividing $2^k$.
Formula \eqref{eq:pyramidpolar} doubles every old coordinate at each step.
We record also the new apex coordinates to check integrality of the whole
polar.

\begin{proposition}\label{prop:family}
The polytope $Q_k$ is reflexive, has dimension $m+k$, and has
\[
 |\vertices(Q_k)|=2^{2^k}\binom{m}{2^k}+k.
\]
\end{proposition}

\begin{proof}
Lemma~\ref{lem:base} gives integral base vertices, and
Lemma~\ref{lem:pyramid} preserves them while adding one integral apex per
step. All $Q_j$ are full-dimensional and contain $0$ in their interiors.
More explicitly, the vertices of $Q_k$ are
\[
 (v,1,\ldots,1)\quad(v\in\vertices(R_{m,r})),\qquad
 (0^m,0^{j-1},-1,1^{k-j})\quad(1\leq j\leq k),
\]
where $0^b$ and $1^b$ mean strings of $b$ zeros and ones, respectively;
a string of length zero is omitted. These points are distinct, giving
the asserted count.

Put $c_i=-2^{k-i}$ for $1\leq i\leq k$.
Iterating \eqref{eq:pyramidpolar} and using $r=2^k$ gives a generating set
of $Q_k^*$ consisting of
\begin{align*}
 &(\pm2^k e_i,c_1,\ldots,c_k) &&(1\leq i\leq m),\\
 &(s,c_1,\ldots,c_k) &&(s\in\{-1,1\}^m),\\
 &b_j=(0^m,0^{j-1},2^{k-j},c_{j+1},\ldots,c_k)
                                      &&(1\leq j\leq k).
\end{align*}
Indeed, each step doubles all old coordinates and appends $-1$.
The polar apex introduced at step $j$ is then doubled $k-j$ times,
which gives $b_j$. This accounts for every generator by induction.
All displayed coordinates are integers, so $Q_k^*$ is a lattice polytope.
Together with the already established properties of $Q_k$, this proves
reflexivity. No intermediate polar is assumed integral.
\end{proof}

\begin{proof}[Proof of Theorem~\ref{thm:main}]
Take $m=40$, $r=32$, $k=5$ and $P=Q_5^*$.
Proposition~\ref{prop:family} shows that $Q_5$ and $P$ are lattice
polytopes. Both have dimension $45$ and contain $0$ in their interiors,
and $P^*=Q_5$. Thus $P$ is reflexive in the standard integer lattice.

For completeness, each vertex $q$ of a full-dimensional polytope $Q$ with
$0\in\inte Q$ gives a facet $\langle q,y\rangle=1$ of $Q^*$.
Separate $q$ strictly from the convex hull of the other vertices to obtain
a functional $c$ uniquely maximized at $q$. Its support value is positive,
so $y=c/h_Q(c)$ satisfies this one polar inequality with equality and all
other vertex inequalities strictly. Small perturbations within that
hyperplane remain feasible, proving that the face has dimension
$\dim Q-1$. Conversely, take $\dim Q$ affinely independent points of a
facet of $Q^*$ and average them. At least one vertex inequality is active
at the average; averaging shows it is active at all the chosen points.
Its hyperplane is the facet's affine hull. Distinct vertices give distinct
hyperplanes, so this correspondence is a bijection.

Consequently,
\[
 F:=f_{44}(P)=|\vertices(Q_5)|
   =2^{32}\binom{40}{32}+5
   =2^{32}\binom{40}{8}+5
   =330303106984181765.
\]
The calculation uses $\binom{40}{8}=76904685$ and $2^{32}=4294967296$.
The comparison can be made entirely with integers:
\begin{align*}
 F^2&=109100142483403824664737926558515225,\\
 6^{45}&=103945637534048876111514866313854976,\\
 F^2-6^{45}&=5154504949354948553223060244660249>0.
\end{align*}
Both quantities being compared are positive, so $F>6^{45/2}$.
This proves the theorem.
\end{proof}

\begin{thebibliography}{3}
\bibitem{HM}
Christian Haase and Ilarion V. Melnikov,
\emph{The reflexive dimension of a lattice polytope},
arXiv:math/0406485v1 (2004).
\href{https://arxiv.org/abs/math/0406485v1}{arXiv:math/0406485v1}.

\bibitem{MMO}
Aki Mori, Kenta Mori and Hidefumi Ohsugi,
\emph{Facet Numbers of Non-Centrally Symmetric Reflexive Polytopes Arising
from Posets}, Discrete \& Computational Geometry (2026).
\href{https://doi.org/10.1007/s00454-026-00879-x}{doi:10.1007/s00454-026-00879-x}.

\bibitem{Nill}
Benjamin Nill, \emph{Gorenstein toric Fano varieties},
arXiv:math/0405448v1 (2004).
\href{https://arxiv.org/abs/math/0405448v1}{arXiv:math/0405448v1}.
\end{thebibliography}

\paragraph{AI use.}
OpenAI Codex assisted the mathematical exploration and proof development,
manuscript drafting, literature checks, and exact-arithmetic checks.
\end{document}
